% ---- pg-371 （书页 359） ----

\begin{mvfigure}
\begin{tikzpicture}[x=0.44cm,y=0.42cm,font=\small,line width=0.5pt]
\draw (0,-1.0) rectangle (16.6,8.6);
\foreach \y in {0,2,4,6,8}{
  \draw (0,\y) -- (0.32,\y);
  \node[left,font=\small] at (-0.05,\y) {\y};}
\foreach \x in {0,5,10,15}{
  \draw (\x,-1.0) -- (\x,-1.28);
  \node[below,font=\small] at (\x,-1.28) {\x};}
\draw[dashed] (0,0) -- (16.6,0);
\draw (0,7.81) -- (16.6,7.81);
\node[right,font=\small] at (16.7,7.81) {UCL};
\node[below,font=\small] at (8.3,-2.05) {时间段};
\node[rotate=90,font=\small] at (-1.35,3.8) {$T^2$};
\foreach \x/\y in {1/0.891,2/4.643,3/3.263,4/2.004,5/2.336,6/2.071,7/2.989,8/3.427,9/3.792,10/0.540,11/0.600,12/8.185,13/7.721,14/1.566,15/0.267,16/0.705}{
  \fill (\x,\y) circle (1.3pt);}
\draw (1,0.891)--(2,4.643)--(3,3.263)--(4,2.004)--(5,2.336)--(6,2.071)--(7,2.989)--(8,3.427)--(9,3.792)--(10,0.540)--(11,0.600)--(12,8.185)--(13,7.721)--(14,1.566)--(15,0.267)--(16,0.705);
\end{tikzpicture}
\phantomsection\label{fig:8.8}\mvcap{图 8.8\quad 基于加班时数的最后三个主成分的 $T^2$ 图}
\end{mvfigure}

\mvsub{控制未来值}

前面我们用将前两个主成分和后 $p-2$ 个主成分分开考虑的办法, 思考过检查一个给定的多元观测值序列是否稳定. 因为 $\chi^2$ 分布曾用于逼近 $T^2$ 图的 UCL 和逼近椭圆形式图的判别距离, 对于监控未来值, 不再需要做任何进一步的修改.

\noindent\textbf{例 8.12}\quad (对于未来主成分的控制椭圆)

例 8.10 中, 我们确定情况 11 是失控的. 我们剔除此点并根据余下的 15 个观测值的协方差, 重新计算特征值和特征向量. 结果示于\hyperref[tab:8.3]{表 8.3} 中.

\begin{mvtable}
\phantomsection\label{tab:8.3}\mvtabcap{表 8.3\quad 来自 15 个稳定观测值的特征向量和特征值}
\begin{tabular}{lccccc}
\toprule
& $\hat{\mathbf{e}}_1$ & $\hat{\mathbf{e}}_2$ & $\hat{\mathbf{e}}_3$ & $\hat{\mathbf{e}}_4$ & $\hat{\mathbf{e}}_5$\\
\midrule
法定加班出勤时数$(x_1)$ & 0.049 & 0.629 & 0.304 & 0.479 & 0.530\\
特殊事件加班时数$(x_2)$ & 0.007 & $-0.078$ & 0.939 & $-0.260$ & $-0.212$\\
处理遗留人员加班时数$(x_3)$ & $-0.662$ & 0.582 & $-0.089$ & $-0.158$ & $-0.437$\\
允许补偿的加班时数$(x_4)$ & 0.731 & 0.503 & $-0.123$ & $-0.336$ & $-0.291$\\
会议时数$(x_5)$ & $-0.159$ & 0.081 & $-0.058$ & $-0.752$ & 0.632\\
\midrule
$\hat{\lambda}_i$ & 2 964 749.9 & 672 995.1 & 396 596.5 & 194 401.0 & 92 760.3\\
\bottomrule
\end{tabular}
\end{mvtable}

主成分已经变了. 主要由特殊事件加班组成的成分, 现在是第三个主成分, 并且不包含在前两个主成分的图中. 因为我们最初的样本大小只是 16, 剔除其一可能造成实质上的差别. 为了设置未来的控制限, 对于一个过程的稳定操作, 至少需要 50 个或更多的观测值.

对加班时间的新前两个主成分的未来数值对, 由\hyperref[fig:8.9]{图 8.9} 给出其 99\% 的预测[见\hyperref[eq:8-36]{式(8-36)}]椭圆. 15 个稳定的主成分对也示于其中. \hfill$\blacksquare$

\begin{mvfigure}
\mvinc{pg-371}{0.564,0.648,0.868,0.841}
\phantomsection\label{fig:8.9}\mvcap{图 8.9\quad 未来加班时间的前两个主成分的 99\%椭圆形式图}
\end{mvfigure}

在化学和制药行业的多元控制的某些应用中, 随时监控的变量超过 100. 这些变量既包括质量变量也包括数字化过程变量. 典型情况是, 与前几个主成分正交的空间, 有大于 100 的维数, 而某
